\section{Related Work}

We situate our contribution across three areas of prior work.

\subsection{Adversarial Examples and Prompt Injection}

Adversarial examples exploit model vulnerabilities through carefully crafted inputs that induce misclassification \cite{goodfellow2014explaining}. In the language model setting, prompt injection attacks manipulate model behavior by embedding adversarial instructions within user inputs \cite{perez2022ignore}. While both phenomena involve models producing outputs misaligned with user expectations, the threat model differs fundamentally from the deception scenario we address. Adversarial attacks originate from an external adversary manipulating inputs; in our setting, the model itself is the source of deceptive outputs, having been instructed (or potentially having learned) to maintain false claims under interrogation. Our interrogator thus targets internally generated deception rather than externally induced failures.

\subsection{Representation Engineering and Probing Classifiers}

Recent work has demonstrated that deception-related information is encoded in model activations and can be extracted through white-box methods. Representation engineering identifies linear directions in activation space corresponding to concepts such as honesty and deception \cite{zou2023representation}. Probing classifiers distinguish truthful from deceptive model states without labeled behavioral data \cite{burns2022discovering}. Azaria and Mitchell~\cite{azaria2023internal} train a classifier on internal activations (SAPLMA) to detect LLM falsehoods; Li et al.~\cite{li2024inference} probe and steer truthfulness-related directions via Inference-Time Intervention (ITI). Both require weight access unavailable for API-only models. Our framework operates under strictly black-box constraints, relying solely on textual exchanges, at the cost of noisier behavioral signals that motivate our multi-turn adaptive strategy. We further note that instruction-following confounds documented in the sycophancy literature \cite{sharma2023sycophancy} directly motivate our prompt-equalized control, which isolates knowledge-state effects from instruction artifacts.

\subsection{LLM Red-Teaming and Evaluation}

Systematic red-teaming has emerged as a key methodology for evaluating language model safety. Large-scale human red-teaming efforts have catalogued failure modes across harmful content categories \cite{ganguli2022red}, while automated red-teaming uses language models themselves to generate adversarial test cases \cite{perez2022red}. However, existing red-teaming paradigms focus primarily on \emph{eliciting} harmful or undesirable outputs---toxicity, bias, dangerous information---rather than \emph{detecting} whether a model's ostensibly normal outputs are deceptive. Our framework addresses this complementary problem: given a model that produces fluent, superficially plausible responses, can an automated interrogator determine whether those responses are truthful?

\paragraph{Summary of the gap.} No prior work combines adaptive multi-turn behavioral interrogation with confidence-based early stopping for black-box instructed lie detection. Representation engineering and probing methods require model internals; red-teaming focuses on eliciting harmful outputs. Our framework operates under strictly black-box constraints, and our central contribution is an empirical characterization of where this approach succeeds and fails.